% TOP_PROBLEM: {"id": 3254, "title": "Cycles in Graphs of Large Chromatic Number", "category_id": 3, "category": {"id": 3, "name": "graph_theory", "display_name": "Graph Theory", "description": "Problems involving graphs, networks, and their properties.", "slug": "graph-theory", "order_index": 3, "created_at": "2026-07-31T15:26:25.671Z"}, "status": "open", "problem_number": "OPG-60001", "set_id": 12, "statusQualification": "The missing hypothesis k>=3 is restored from the published formulation and explicitly explained."}
% FIRST_POSTED: 2026-10-01
% ENTRY_KIND: statement_only
% UnsolvedMath ID 3254; source catalogue code OPG-60001
% !TeX program = lualatex
% Definition and sources only; this file is not an LLM solution attempt.
\documentclass[11pt,a4paper]{article}
\usepackage[margin=25mm]{geometry}
\usepackage{fontspec}
\setmainfont{lmroman10-regular.otf}[BoldFont=lmroman10-bold.otf,
 ItalicFont=lmroman10-italic.otf,BoldItalicFont=lmroman10-bolditalic.otf]
\usepackage{amsmath,amssymb,mathtools}
\usepackage{unicode-math}
\setmathfont{latinmodern-math.otf}
\AtBeginDocument{\let\setminus\smallsetminus}
\usepackage{xurl,hyperref}
\hypersetup{colorlinks=true,urlcolor=blue}
\setcounter{secnumdepth}{0}
\setlength{\parindent}{0pt}
\setlength{\parskip}{5pt}
\setlength{\emergencystretch}{3em}
\begin{document}
\section*{Cycles in Graphs of Large Chromatic Number}\label{3254}
{\small UnsolvedMath ID 3254 · OPG-60001 · Graph Theory · OpenGarden}
\subsection{Definitions and mathematical statement}
Let $G$ be a finite simple undirected graph and let $k\ge3$ be an integer.
The chromatic number $\chi(G)$ is the least number of colours in a proper
vertex colouring. A cycle means a simple cycle, with no repeated vertex
except the equality of its initial and final vertex. Its length is its
number of edges. Count cycles as distinct subgraphs, without distinguishing
a starting vertex or a direction of traversal; chords between vertices
of a cycle are permitted.

Let $c_{0,k}(G)$ count the cycles whose lengths are positive multiples of
$k$. The intended conjecture is
\[
 \chi(G)>k\quad\Longrightarrow\quad
 c_{0,k}(G)\ge\frac{(k+1)(k-1)!}{2}.
\]
The bound counts cycles individually, with no requirement of edge or
vertex disjointness. In $K_{k+1}$ the only possible positive multiple of
$k$ among cycle lengths is $k$. Choosing the omitted vertex and a cyclic
ordering of the others gives exactly $(k+1)(k-1)!/2$ such cycles, explaining
the proposed extremal quantity.

The restriction $k\ge3$ is essential and is stated explicitly in the
primary formulation [S3]. It is missing from the imported one-sentence
statement. Without it the assertion fails: $K_3$ is more than
2-chromatic but has no even cycle, and a nontrivial tree is more than
1-chromatic but has no cycle. The present expansion restores this
published parameter range rather than treating those trivial failures
as the intended problem.
\subsection{Short English statement}
Does every graph needing more than k colours contain at least as many cycles of length divisible by k as the complete graph on k+1 vertices, for k at least three?
\subsection{Sources}
\begin{itemize}
\item[S1] ULAM.AI and the UnsolvedMath Contributors, supplied problems.json, record 3254 (OPG-60001), OpenGarden collection. Catalogue identity and source transcription; CC BY 4.0.
\url{https://huggingface.co/datasets/ulamai/UnsolvedMath}.
\item[S2] Open Problem Garden, Cycles in Graphs of Large Chromatic Number. Original problem and discussion; the mathematical formulation below is expanded from the supplied source record.
\url{http://www.openproblemgarden.org/op/cycles_in_graphs_of_large_chromatic_number}.
\item[S3] S. Kim and M. E. Picollelli, 4-Chromatic Graphs Have At Least 4 Cycles of Length 0 mod 3 (2023), abstract explicitly states k at least 3.
\url{https://arxiv.org/abs/2312.05945}.
\end{itemize}
\end{document}
